\documentclass{article}
\usepackage{enumitem}
\usepackage{amsmath}
\begin{document}

\title{Chapter 4: Animal Kingdom}
\date{}
\maketitle

\begin{enumerate}[label=Q\arabic*.]

\item Which of the following criteria is used as the basis for classifying animals into diploblastic and triploblastic?
\\(A) Number of appendages
\\(B) Number of embryonic germ layers
\\(C) Type of body cavity
\\(D) Mode of reproduction

\item An animal that has a true coelom lined by mesoderm on both sides is called:
\\(A) Acoelomate
\\(B) Pseudocoelomate
\\(C) Coelomate
\\(D) Haemocoelomate

\item Which phylum is characterized by presence of cnidoblasts (nematocysts)?
\\(A) Porifera
\\(B) Cnidaria
\\(C) Platyhelminthes
\\(D) Echinodermata

\item The water vascular system is a characteristic feature of:
\\(A) Arthropoda
\\(B) Mollusca
\\(C) Echinodermata
\\(D) Hemichordata

\item Which of the following correctly describes the characteristic of Porifera?
\\(A) They have true tissues and organs
\\(B) They are diploblastic with canal system and choanocytes
\\(C) They show radial symmetry exclusively
\\(D) They reproduce only by budding

\item Assertion: Nematoda are pseudocoelomates.\\
Reason: Their body cavity is not lined by mesoderm on both sides.
\\(A) Both true, Reason is correct explanation
\\(B) Both true, Reason is not correct explanation
\\(C) Assertion true, Reason false
\\(D) Assertion false, Reason true

\item Which feature distinguishes Annelida from Arthropoda?
\\(A) Presence of segmentation
\\(B) Presence of jointed appendages in Arthropoda
\\(C) Presence of coelom
\\(D) Bilateral symmetry

\item The open circulatory system is found in:
\\(A) Annelida and Echinodermata
\\(B) Arthropoda and Mollusca
\\(C) Chordata and Annelida
\\(D) Platyhelminthes and Nematoda

\item Which of the following is a correct match?
\\(A) \textit{Balanoglossus} -- Echinodermata
\\(B) \textit{Amphioxus} -- Urochordata
\\(C) \textit{Herdmania} -- Urochordata
\\(D) \textit{Petromyzon} -- Chondrichthyes

\item Which characteristic is unique to chordates and not found in any non-chordate phylum?
\\(A) Bilateral symmetry
\\(B) Post-anal tail at some stage of life
\\(C) Segmentation
\\(D) Coelom

\item In which class of vertebrates is the heart two-chambered?
\\(A) Amphibia
\\(B) Reptilia
\\(C) Pisces (fish)
\\(D) Aves

\item Which of the following animals is oviparous and belongs to Class Mammalia?
\\(A) \textit{Ornithorhynchus} (Platypus)
\\(B) \textit{Macropus} (Kangaroo)
\\(C) \textit{Delphinus} (Dolphin)
\\(D) \textit{Pteropus} (Flying fox)

\item Match the following phyla with their excretory organs:\\
1. Platyhelminthes $\rightarrow$ (a) Flame cells\\
2. Annelida $\rightarrow$ (b) Nephridia\\
3. Arthropoda $\rightarrow$ (c) Malpighian tubules\\
4. Mollusca $\rightarrow$ (d) Kidneys (metanephridia)
\\(A) 1-a, 2-b, 3-c, 4-d
\\(B) 1-b, 2-a, 3-d, 4-c
\\(C) 1-c, 2-d, 3-a, 4-b
\\(D) 1-a, 2-d, 3-b, 4-c

\item Metameric segmentation is most prominently seen in:
\\(A) Platyhelminthes
\\(B) Annelida
\\(C) Mollusca
\\(D) Echinodermata

\item Which of the following is a cold-blooded (poikilothermic) vertebrate?
\\(A) Pigeon
\\(B) Bat
\\(C) Crocodile
\\(D) Whale

\end{enumerate}
\end{document}
